\section{Step 66 Bound each expansion before storing its bytes}\label{png-huff-s66}
Fill \texttt{huff\_store\_lengths}. Reload current cursor and combined end, compare the prepared repetition count with remaining entries, and reject an overrun before any store. Write the prepared length byte that many times, advance the cursor and retain its new value. The supplied continuation reads another operation unless the range is complete.

The first direct operation writes length one at combined index zero and moves the cursor to index one. The same region will later handle repeat runs. The literal/length and distance portions form one sequence, so a permitted repeat may cross their boundary; the combined end, not that boundary, limits expansion.

\paragraph{Observe} Use \texttt{s66-1}; stop before \texttt{huff\_lengths\_stored}. Inspect the state described above before continuing.

\Needspace{7\baselineskip}
\paragraph{Check} Which boundary limits a repeat expansion?
\begin{itemize}
  \item[A.] The first alphabet boundary even when more declared entries remain.
  \item[B.] The current code’s bit length.
  \item[C.] The end of the declared combined length vector.
\end{itemize}
